% !TEX root = ../main.tex
\section{Tile Size of Gathered Weights}
\label{app:tiles}

\begin{table}[H]
\caption{Rank sort on 12-layer cuts, measured before the attention fix: latency change against the token-owned program
at capacity $T$ (default tiles) for three maximum tile sizes.}
\label{tab:tiles}
\vskip 0.1in
\begin{center}
\begin{small}
\begin{tabular}{@{}rrrr@{}}
\toprule
$T$ & Default tiles & 512\,KiB & 1024\,KiB \\
\midrule
128 & $+4.7\%$ & $\mathbf{-1.4\%}$ & -- \\
256 & $-5.1\%$ & $\mathbf{-18.2\%}$ & $-7.9\%$ \\
512 & $-12.7\%$ & $-9.3\%$ & $\mathbf{-15.9\%}$ \\
\bottomrule
\end{tabular}
\end{small}
\end{center}
\vskip -0.15in
\end{table}

\textbf{Tile size.} Table~\ref{tab:tiles} shows the sensitivity of gathered weights to the compiler's maximum tile size.
With default tiles the rank sort is 4.7\% slower than the token-owned program at capacity $T$ at $T=128$; with 512\,KiB
tiles it is 1.4\% faster. The
best size moves with the chunk, 512\,KiB up to 256 tokens and 1024\,KiB at 512; static weights prefer the default only
at $T=128$ (small tiles slow them by 1--2\%), while at 256 and 512 tokens 512\,KiB tiles make the static program 8\% and 16\%
faster (Table~\ref{tab:attr}). Each program therefore uses its own best size. The tile size is the granularity of the
compiler's row skipping (Appendix~\ref{app:attrprof}); at $T=128$ a lane is one tile at every size.
